\documentclass[twocolumn]{article}
\usepackage[margin=.7in]{geometry}
\usepackage[showisoZ=false]{datetime2}
\usepackage{graphicx}
\usepackage{url}
\usepackage{tabularx}
\usepackage{amsthm}
\usepackage{amsmath}
\usepackage{amssymb}
\newtheorem{security}{security interpretation}
\newtheorem{definition}{definition}
\newtheorem{lemma}{lemma}
\newtheorem{theorem}{theorem}
\newtheorem{discovery}{discovery}
\newtheorem{law}{law}
\newtheorem{advertisement}{advertisement}
\usepackage[linesnumbered,ruled,vlined]{algorithm2e}
\DeclareMathOperator{\fexists}{\mathtt{exists}}
\DeclareMathOperator{\fread}{\mathtt{read}}
\DeclareMathOperator{\fwrite}{\mathtt{write}}
\DeclareMathOperator{\sizeof}{\mathtt{sizeof}}
\DeclareMathOperator{\delete}{\mathtt{delete}}
\DeclareMathOperator{\enc}{\mathtt{enc}}
\DeclareMathOperator{\dec}{\mathtt{dec}}
\DeclareMathOperator{\maxf}{max}
\DeclareMathOperator{\hash}{\mathtt{hash}}
\DeclareMathOperator{\rhash}{\mathtt{rhash}}
\DeclareMathOperator{\mhash}{\mathtt{mhash}}
\DeclareMathOperator{\argon}{\mathtt{argon2}}
\DeclareMathOperator{\ciphart}{\mathtt{ciphart}}
\DeclareMathOperator{\cost}{\mathtt{cost}}
\renewcommand{\contentsname}{paper's layout}
\DTMsetdatestyle{iso}
\usepackage{cleveref} % must be loaded last
\begin{document}
\SetAlgorithmName{algorithm}{}{list of algorithms}
\SetInd{.15em}{1em}

\begin{center}
\Huge
ciphart\\
\Large
faster memory-\emph{harder} key derivation \\
with easier security interpretation\\
\normalsize
\emph{caveman}\footnote{mail: \texttt{toraboracaveman [at] protonmail [dot]
com}.}\\
\footnotesize
\DTMnow\\
\end{center}

\noindent\textbf{synopsis---}\emph{argon2}\footnote{\url{https://github.com/P-H-C/phc-winner-argon2}}
is a fast and simple memory-hard key derivation function.  compared to
\emph{scrypt}\footnote{\url{http://www.tarsnap.com/scrypt/scrypt.pdf}},
\emph{argon2} is better.  but i claim that \emph{argon2} is not fast
enough, not memory-hard enough, and its contribution to our security is not
simple enough to understand.

henceforth, i propose \emph{ciphart}, which is:
\begin{itemize}
    \item easier --- because its security contribution is measured in the
    unit of shannon's entropy.  i.e. when \emph{ciphart} derives a key for
    you, it tells you that it has \emph{injected} a specific guaranteed
    quantity of shannon's entropy bits into your derived key.  this is
    possible thanks to my invention, the ``entropy injection'' theorem.

    this offers a great help as it gives us yet another much simpler
    approach to quantify our security gain as opposed to being limited to
    surveying the industry of application-specific integrated circuits as
    done in the \emph{scrypt} paper.

    \item harder --- because it can require crazy-large amounts of memory,
    beyond our random-access memory, thanks to it being able to use the
    hard-disk as well.  this is possible thanks to my discovery ``cacheable
    keys''.

    this is optional, but i extremely like it as it effectively gives me
    much more security while eventually becoming much faster as well, and
    the adversary cannot get my cache even if he steals my hard-disks.

    \item faster --- because it does not abuse hashing functions.  it uses
    hashing functions when using them is more suited, and uses symmetric
    block encryption functions when using them is  more suited.  this is
    thanks to my ``hashing is only for compression'' law.

    \emph{argon2} incorrectly limits itself to only use a hashing function.
    at the surface it may appear simpler, but it is actually more complex
    as it ends up re-inventing what resembles a symmetric block encryption
    function off the hashing function, except for being slower and with
    needless potential entropy loss.
\end{itemize}

\texttt{libciphart}\footnote{\url{https://github.com/Al-Caveman/libciphart}}
is a library that implements \emph{ciphart} very closely to this paper,
without much fluff.  this should make integrating \emph{ciphart} into other
systems more convenient.

\texttt{ciphart}\footnote{\url{https://github.com/Al-Caveman/ciphart}} is
an application for encrypting and decrypting files that makes use of
\texttt{libciphart}.  this application is intended for use by end-users or
scripts, henceforth it has some fluff to treat mankind with dignity.

\break
\tableofcontents

\section{background}
\subsection{passwords' entropy}
\begin{definition}[specifications set]
    the specifications set of a thing $x$ is a set $\mathcal{S}_x$ that
    contains binary specifications, where for any specification $s \in
    \mathcal{S}_x$, $s \in \{0,1\}$, such that $x$ becomes fully known if,
    and only if, the value for every specification in $\mathcal{S}_x$ is
    known.
\end{definition}

say that $\mathcal{S}_x$ is the specifications set of thing $x$.  generally,
a scalable way to reveal those specifications is to ask questions, such
that if each one of them is answered, the quantity of unknown
specifications shrinks in half or more.  let's stick to shrinkage in half;
i.e. balanced binary trees.  this means that the total number of such
balanced binary questions is $\log_2 |\mathcal{S}_x|$, where $|\mathcal{S}_x|$
is the quantity of total specifications in $\mathcal{S}_x$ about $x$.

but is that scalable approach, also the most efficient one in extracting
information about $x$'s specifications?  the answer is: it depends.  the
reason is that a binary search tree is not necessarily the most informative
one in revealing information about the specifications in $\mathcal{S}_x$.
e.g. is it better to reveal, say $10$ specifications that we already sort
of knew their value for a probability of, say, $0.9999$ than to reveal,
say, $5$ specifications that their their values completely blew our minds
with a probability of, say, $0.0001$?

shannon's entropy tells us the minimum number of questions to be asked in
order to extract all information about $x$'s specifications,
$\mathcal{S}_x$:
\begin{equation}
    H(\mathcal{S}_x) = \sum_{s \in \mathcal{S}_x}
        \Pr(s=0) \log_2 \Pr(s=0)^{-1}
\end{equation}

so when is $\log_2 |\mathcal{S}_x|$ optimal?  $H(\mathcal{S}_x) = \log_2
|\mathcal{S}_x|$ if every specification in $\mathcal{S}_x$ is equally
informative to us, which is the case when there is no redundancy in the
specifications.  i.e.  for any $s \in \mathcal{S}_x$, $\Pr(s=0) =
|\mathcal{S}_x|^{-1}$.
\begin{proof}
\begin{equation}
\begin{split}
    H(\mathcal{S}_x)
    &= \sum_{s \in \mathcal{S}_x} |\mathcal{S}_x|^{-1} \log_2
        (|\mathcal{S}_x|^{-1})^{-1} \\
    &= \sum_{s \in \mathcal{S}_x} |\mathcal{S}_x|^{-1} \log_2
        |\mathcal{S}_x| \\
    &= |\mathcal{S}_x| |\mathcal{S}_x|^{-1} \log_2 |\mathcal{S}_x| \\
    &= \log_2 |\mathcal{S}_x| \\
\end{split}
\end{equation}
\end{proof}

in other words, $H(\mathcal{S}_x) = \log_2 |\mathcal{S}_x|$ if
$\mathcal{S}_x$ is optimally compressed with no redundancy at all.  say
that $\mathcal{S}_x^*$ is the same as $\mathcal{S}_x$, except that the
former is optimally compressed to be totally free of any redundancy.  in
this case, it must follow that $H(\mathcal{S}_x^*) = H(\mathcal{S}_x) =
\log_2 |\mathcal{S}_x^*|$.

let's say that $p$ is a password with an entropy of, say, $20$ bits.  what
this means is that $H(\mathcal{S}_p) = H(\mathcal{S}_p^*) = \log_2
|\mathcal{S}_p^*| = 20$.  also let's say that the adversary has an $8V$
bits hash $k \gets \hash(p, 8V)$ of the password $p$, and that he wants to
find $p$ that gave $k$.

when we talk about brute-forcing to find password $p$, we talk about
performing the most efficient brute-forcing, which is brute-forcing against
the redundancy-free specifications set $\mathcal{S}_p^*$.  we certainly do
not talk about brute-forcing $\mathcal{S}_x$, because it contains
redundancies in its specifications, which is needlessly wasteful.

this is also why we say that the total number of password candidates to
test is $2^{H(\mathcal{S}_p)} = |\mathcal{S}_p^*|$, because this is the
number of binary specifications in the optimally compressed specifications
set $\mathcal{S}_p^*$.  for any $s^* \in \mathcal{S}_p^*$, $s^*$ looks
something like ``is $\hash(\hat p, 8V) = k$?'', where $\hat p$ is a
candidate password.  there are $2^{H(\mathcal{S}_p)}$ many such password
candidates, each represented in a binary question in $\mathcal{S}_p^*$.

\subsection{passwords' security}
\begin{definition}[systems' security]
the security of a system is the cost of the cheapest method that can break
it.
\end{definition}

\subsubsection{$\hash$}
say that we've got a $8V$ bit key $k \gets \hash(p, 8V)$, derived from
password $p$.  say that the attacker has $k$ but wants to figure out $p$.

in this case, he will need to brute-force the password space in order to
find $p$ that gives $k$.  in order to make his brute-forcing cheapest
possible, he will need to brute-force the redundancy-free specifications
set $\mathcal{S}_p^*$, which contains $2^{H(\mathcal{S}_p)}$ many binary
tests, each of which asks whether a password candidate $\hat p$ gives given
a key candidate $\hat k$ that matches $k$.  his cost is:
\begin{equation}\label{eq_cost_passbruteforce}
    2^{H(\mathcal{S}_p)} \left(
        \cost(\hash) + \cost(\text{if } \hat k = k)
    \right)
\end{equation}

one way to estimate $\cost$ is to survey the asics industry.  by surveying
the asics industry to get an idea how much money it costs to get a given
key, or password, space brute-forced within a target time
frame\footnote{see the \emph{scrypt} paper for an example.}.  this has an
expensive housekeeping and is usually not possible to get any guarantees as
we don't know about state-of-art manufacturing secrets that adversaries may
have.

another way is to ignore anything that has no cryptographic guarantee.  so,
in (\ref{eq_cost_passbruteforce}), cryptography
guarantees\footnote{statistically by confidence earned through peer review
and attempts to break encryption algorithms.} that $2^{H(\mathcal{S}_p)}$
many $\hash$ calls are performed and that many equality tests.  the $\hash$
call needs to be done once, so let's give it a unit of time $1$.  the
equality test also needs to be called once, but since since it's so cheap
it's easier to just assume that its cost is free.  this way
(\ref{eq_cost_passbruteforce})
becomes just:
\begin{equation}\label{eq_simplecost_passbruteforce}
    2^{H(\mathcal{S}_p)} (1+0) = 2^{H(\mathcal{S}_p)}
\end{equation}

further, for convenience i guess, it seems that people report it in the
$\log_2$ scale.  i.e.:
\begin{equation}\label{eq_pass_entropy}
    \log_2 2^{H(\mathcal{S}_p)} = H(\mathcal{S}_p)
\end{equation}

i think this is why people use shannon's entropy of passwords as a measure
of their security.  not because it is the quantity of security, but rather
because its the quantity of \emph{simplified} security.

i like using shannon's entropy as a measure of simplified security
quantity, so i'm going to build on it.

\subsubsection{recursive $\hash$: $\rhash$}
if the $\hash$ function is replaced by an $N$-deep recursion over $\hash$,
like:
\[
    \begin{split}
        & \rhash(p, 8V, N) \\
    ={} &  \hash(\hash(\ldots\hash(p, 8V), \ldots, 8V), 8V)
    \end{split}
\]
then, if $\hash$ is not broken,  (\ref{eq_cost_passbruteforce}) becomes:
\begin{equation}\label{eq_cost_passbruteforce_N}
    2^{H(\mathcal{S}_p)} \left(
        N\cost(\hash) + \cost(\text{if } \hat k = k)
    \right)
\end{equation}
(\ref{eq_simplecost_passbruteforce}) becomes:
\begin{equation}\label{eq_simplecost_passbruteforce_N}
    \begin{split}
    2^{H(\mathcal{S}_p)} (N+0) &= N2^{H(\mathcal{S}_p)} \\
                  &= 2^{H(\mathcal{S}_p) + \log_2 N}
    \end{split}
\end{equation}
and the $\log_2$ scaled version becomes:
\begin{equation}\label{eq_pass_cavemanentropy_rhash}
    H(\mathcal{S}_p) + \log_2 N
\end{equation}

\subsubsection{memory-hard $\hash$: $\mhash$}
let $\mhash$ be like $\rhash$, except that it also requires $M$ many memory
bytes such that, as available memory is linearly reduced from $M$, penalty
in cpu time grows exponentially.  let $M$ be requested memory, $A$ be
available memory, and $e(M - A)$ be the exponential penalty value for
reduction in memory, where $e(0) = 1$.
\begin{equation}
    \begin{split}
        & \cost\Big(\mhash(p, N, M)\Big) \\
    ={} & \cost\Big(\rhash(p, N)\Big)^{e(M-A)}
    \end{split}
\end{equation}

if $\hash$ in (\ref{eq_cost_passbruteforce}) is replaced by the $M$-bytes
memory-hardened $N$-deep recursion hash function $\mhash$, then
(\ref{eq_cost_passbruteforce}) becomes:
\begin{equation}\label{eq_cost_passbruteforce_NM}
    2^{H(\mathcal{S}_p)} \left(
        N^{e(M-A)}\cost(\hash) + \cost(\text{if } \hat k = k)
    \right)
\end{equation}
(\ref{eq_simplecost_passbruteforce}) becomes:
\begin{equation}\label{eq_simplecost_passbruteforce_NM}
    \begin{split}
    2^{H(\mathcal{S}_p)} (N^{e(M-A)}+0) &= N^{e(M-A)} 2^{H(\mathcal{S}_p)} \\
                  &= 2^{H(\mathcal{S}_p) + \log_2 N^{e(M-A)}} \\
                  &= 2^{H(\mathcal{S}_p) + e(M-A)\log_2 N}
    \end{split}
\end{equation}
and $\log_2$ scaled version becomes:
\begin{equation}\label{eq_pass_cavemanentropy_mhash}
    H(\mathcal{S}_p) + e(M-A)\log_2 N
\end{equation}

\subsubsection{caveman's entropy}
\begin{definition}[caveman's entropy]\label{def_cavemanentropy}
it is clear that (\ref{eq_pass_entropy}) is entropy, but i didn't make it
clear yet what (\ref{eq_pass_cavemanentropy_rhash}) and
(\ref{eq_pass_cavemanentropy_mhash}) are, so i will give them a temporary
name until i tell you what they are later on in this paper.  for now, let's
call it \emph{caveman's entropy}, $C$, which i define as follows:
\[
\begin{split}
C\Big(p, \hash(\ldots)\Big) &= H(\mathcal{S}_p) \\
C\Big(p, \rhash(\ldots, N)\Big) &= H(\mathcal{S}_p) + \log_2 N \\
C\Big(p, \mhash(\ldots, N, M)\Big) &= H(\mathcal{S}_p) + e(M-A)\log_2 N \\
\end{split}
\]
\end{definition}

\section{fundamental ideas}
\subsection{``entropy injection'' theorem}
\begin{definition}[bundled specifications set]
    $\mathcal{\hat S}_x = \{\hat s_0, \hat s_1, \ldots\}$ is a bundled
    specifications set of $x$'s specifications set $\mathcal{S}_x = \{s_0,
    s_1, \ldots\}$ if there exists $\hat s_i \in \mathcal{\hat S}_x$ and $j
    \ge 1$ such that $\hat s_i = \{s_{ji+0}, s_{ji+1}, \ldots,
    s_{ji+j-1}\}$.
\end{definition}

\begin{lemma}[specifications bundling is lossless]
    if $\mathcal{\hat S}_x$ is a bundled specifications set of
    $\mathcal{S}_x$, then $H(\mathcal{\hat S}_x) = H(\mathcal{S}_x)$
    because bundling does not reduce number of questions that need to be
    asked in order to reveal information about the specifications, as such
    bundling merely groups certain specifications inside a
    super-specification, and such grouping does not delete information.
\end{lemma}

say that $p$ is a password that all ambiguities about it are resolved after
getting, on average, $H(\mathcal{S}_p)$ many balanced binary questions,
$q_0, q_1, \ldots, q_{H(\mathcal{S}_p)-1}$, answered.

case 1 --- say that the adversary got $p$'s hash, $k \gets \hash(p \Vert s,
8V)$, and that his goal is to find out $p$ by brute-forcing.  in this case,
the adversary will be considering password candidates $p_0, p_, \ldots,
p_{2^{H(\mathcal{S}_p)-1}}$, and tests them by asking the following binary
questions in order to guarantee finding $p$:
\begin{itemize}
    \item $q_0$: is $\hash(p_0 \Vert s, 8V) = k$?
    \item $q_1$: is $\hash(p_1 \Vert s, 8V) = k$?

    $\vdots$

    \item $q_{H(\mathcal{S}_p)-1}$: is $\hash(p_{H(\mathcal{S}_p)-1} \Vert
    s, 8V) = k$?
\end{itemize}

case 2 --- say that $\mhash$,was used instead of just $\hash$;  i.e. the
adversary got $\hat k \gets \rhash(p \Vert s, 8V, N)$ instead.  in this
case, the adversary will seek to answer the following binary questions:
\begin{itemize}
    \item $\hat q_0$: is $\rhash(p_0 \Vert s, 8V, N) = k$?
    \item $\hat q_1$: is $\rhash(p_1 \Vert s, 8V, N) = k$?

    $\vdots$

    \item $\hat q_{H(\mathcal{S}_p)-1}$: is $\rhash(p_{H(\mathcal{S}_p)-1}
    \Vert s, 8V, N) = k$?
\end{itemize}

so, you may say that we're still in need to ask, on average,
$H(\mathcal{S}_p)$ many balanced binary questions, right?  i disagree.

i think what is happening with $\rhash$ is that each question $\hat q_0,
\hat q_1, \ldots \hat q_{H(\mathcal{S}_p)-1}$ is a \emph{bundled} question,
each, containing $N$ many balanced binary questions.

i.e., for any question $\hat q_i$, answering it is equivalent to answering
the $N$-deep recursion over $\hash$:
\[
    \text{$\hat q_i$: is }\hash(\hash(\ldots\hash(p_i \Vert s, 8V), \ldots,
    8V), 8V) = k \text{?}
\]
which its answer is ``yes'' if and only if all of the following $N$ many
questions are answered by ``yes'':
\begin{itemize}
    \item is $\hash(p_i \Vert s, 8V) = k_0$?
    \item is $\hash(k_0 \Vert s, 8V) = k_1$?
    \item is $\hash(k_1 \Vert s, 8V) = k_2$?

    $\vdots$

    \item is $\hash(k_{N-1} \Vert s, 8V) = k$?
\end{itemize}

meaning, if any $\hat q_i$ question is answered, it necessarily means that
we have answered $N$ many $q_i$-like questions.  if answer to $\hat q_i$ is
``yes'', it means that the answer to each of those $N$ many $q_i$-like
questions is also ``yes''.  if answer to $\hat q_i$ is ``no'', it also
means that the answer to each of those $N$ many questions is also
``no''\footnote{unless hashing collisions happen, which is very unlikely.}.

in other words, functions $\rhash$ and $\mhash$ are using cryptography to
force the adversary to need to answer more questions beyond
$H(\mathcal{S}_p)$.  $\rhash$ increases questions by a factor of $N$.
$\mhash$ increases questions by a factor of $N^{e(M-A)}$.

\begin{theorem}[entropy injection] \label{theorem_entropy_injection}
    for any password $p$, and any positive numbers $V$, $N$ and $M$:
    \[
        H\Big(\rhash(p\Vert s, 8V, N)\Big) = C\Big(p, \rhash(\ldots, N)\Big)
    \]
    \[
        H\Big(\mhash(p\Vert s, 8V, N, M)\Big) = C\Big(p, \mhash(\ldots, N,
        M)\Big)
    \]
\end{theorem}

\subsection{``cacheable keys'' discovery}
\begin{discovery}[cacheable keys]\label{discov_key_caching}
    caching keys securely is easily doable, and great security utility
    exists in doing so for expensively-derived keys.
\end{discovery}

\subsubsection{why does it work}
\begin{itemize}
    \item when expensively derived keys are cached, only the first key
    derivation call will be expensive, while subsequent calls will be
    semi-instantaneous.  this effectively allows users to tolerate much
    more expensive, or secure, key derivation as it only happens during the
    initial, say, login phase.  subsequent use of the extremely expensive
    key is instantaneous.

    so instead of having the user use a somewhat expensive key derivation
    by wayting say, $3$ seconds in each login, he will ---instead--- wait,
    say $10$ seconds in his initial login in order to utilise a much more
    expensive key derivation, and then wait near $0$ seconds for every
    subsequent login as the expensive key is cached.

    \item derived keys can be cached securely, without increasing most
    users' assumptions.  e.g. cached keys can live in a \emph{dm-crypt}
    partition that is encrypted with a large encryption key that is stored
    properly, and the cache can have strict read permissions so that only
    the unique user that runs \texttt{ciphart} executable can read it.

    this only requires to trust the user \texttt{root}, which is already
    trusted by almost everyone.  so we are not introducing a new
    assumption.
\end{itemize}

for most people, if \texttt{root} is compromised, then the adversary
can break every other key derivation function, including those that do
not cache keys, by simply, say, running a keylogger.

so, practically, we are not increasing the assumptions, but we are only
increasing the value that we can extract from the assumptions that we
already have.

most importantly, utilising discovery \ref{discov_key_caching} allows us to
achieve a memory-\emph{harder} key derivation in an extremely  usable way.
more details on memory-harderness later.

\subsubsection{potential adversary strategies}
let's see what may the adversary try to do against a key caching system:
\begin{itemize}
    \item adversary strategy 1 --- hack into user's system and execute a
    program as that user that tries to read the password cache file to
    obtain the memory-harder key inside it.

    \textbf{answer:} he will not be able to read the file due to strict
    read permissions of the cached files as set earlier.

    \item adversary strategy 2 --- steal user's hard disk and try to mount
    it in his system, to login as root and chance permissions of the key
    cache files.

    \textbf{answer:} the partition where the cached files are saved are
    properly encrypted, so he can't see the cached files, let alone
    changing their read permissions.

    \item adversary strategy 3 --- break into machine's \texttt{root}
    account.

    \textbf{answer:} he will succeed, but then he can also run a keylogger,
    which will also break every other key derivation function, even those
    who do not use key caching, such as \emph{scrypt}, \emph{argon2}, etc.
\end{itemize}

\subsection{``hashing is only for compression'' law}
which one is simpler when, say, building a wooden house?
\begin{itemize}
    \item option 1 --- use only nails and, when you need screws,
    modify some nails into screws.
    \item option 2 --- use nails and screws.
\end{itemize}

on the surface, option 1 may appear as the simpler choice as it only uses
nails, while option 2 uses both nails ans crews.

but a deeper look shows that option 1 is actually a lie, as it is also
using screws alongside nails, except that the screws are constrained by
being re-invented by modifying nails.  in other words, option 1 has the
extra assumption that its screws must be made using nails, while option 2
does not have this extra assumption.  hence option 2 is actually simpler.

my answer with \emph{ciphart} is that option 2 is the simpler choice
because it removes the re-invention aspect, specially if the re-invented
screws were worse than screws that were made as screws from the start.

on the other hand, when \emph{argon2} acts as if option 1 is better, which
is wrong at every level.

before i write about how \emph{argon2} is an example of adopting the
mistake in option 1, i want to define the nails and the screws of a key
derivation function, strictly from the perspective of key derivation
functions.

\begin{definition}[hashing functions as seen by key derivation
functions]\label{def_kdf_hash}
    a function that maps input to output such that:
    \begin{itemize}
        \item unlimited input --- input is an unbounded number of
        concatenated data chunks.
        \item compression --- input's shannon's entropy bits could be
        larger than the total bit size of the output.
        \item preserving entropy is tried --- output should have as much
        entropy bits from the input, but entropy loss is possible.
        \item walking backwards is extremely hard --- analysing the output
        to find the input is computationally too hard.
    \end{itemize}
\end{definition}

\begin{definition}[symmetric block encryption functions as seen by key
derivation functions]\label{def_kdf_enc}
    a function that maps input to output such that:
    \begin{itemize}
        \item limited input --- fixed sized data and a fixed sized
        key.
        \item preserving entropy is guaranteed --- no entropy loss is
        possible.  the proof is that, if the encryption key is known, we
        can bring back every input bit from the output by decryption.
        \item walking backwards is extremely hard --- analysing the output
        to find the input is computationally too hard.
    \end{itemize}
\end{definition}

when \emph{argon2} completes solving the tasks in its memory pad, it
derives the output key by hashing certain chunks of bytes in the pad.  the
number of hashed chunks depends on pad's size, so it's not of a fixed size,
and henceforth meets the properties of a hashing function shown in
\cref{def_kdf_hash}.  therefore, \emph{argon2} using a hashing function at
this stage is justified.

when \emph{argon2} is solving tasks in the memory pad, it still uses a
hashing function, despite the fact that it is strictly dealing with inputs
of fixed lengths, which also meets the properties of a symmetric block
encryption function in \cref{def_kdf_enc}.  here, \emph{argon2} better use
a symmetric block encryption function instead of a hashing function.  using
a hashing function here is a problem due to:
\begin{itemize}
    \item re-invention of wheels --- having \emph{argon2} concatenate two
    inputs of a fixed size in order to derive an output of the same size is
    effectively an attempt to re-invent a symmetric block encryption
    function.  i.e. the concatenation is used to emulate the effect of
    having a pre-shared key which exists in symmetric block encryption
    functions.  why re-invent keys by concatenation when there exists
    functions that already have keys?

    \item needless risk of entropy loss --- using a hashing function when
    dealing with fixed input sizes needlessly increases the probability of
    having potential entropy losses.  this is due to the fact that hashing
    functions \emph{try} to preserve input's entropy, but cannot guarantee
    it, while symmetric block ciphers \emph{do} guarantee it.  so why have
    the possibility of losing entropy bits when you don't have to?

    \item slower memory filling rate --- generally speaking, hashing
    functions tend to be slower than symmetric block encryption functions.
    this is because of dealing with compression is harder than not.

    symmetric block ciphers guarantee preserving input's entropy in the
    output with much less effort thanks to the fact that the output is at
    least as large as the input.

    but hashing functions don't have the luxury of having an output that's
    as large as the input, thus they need to work a lot more in order to
    ensure that no input entropy is needlessly lost.

    this slowness is bad as it reduces number of passes over the memory pad in
    a unit of time.  more passes over the memory pad are important for
    strengthening the memory hardness.
\end{itemize}

\begin{law}[hashing is only for compression]\label{law_simplification}
    use hashing functions only when compression happens.  otherwise, use
    symmetric block encryption functions.
\end{law}

\section{ciphart}
\subsection{the algorithm}
\subsubsection{parameters}
\begin{tabularx}{\columnwidth}{lX}
    $P$ & password.\\
    $S$ & salt.\\
    $M$ & total random-access memory in bytes.\\
    $D$ & total hard-disk memory in bytes.\\
    $F$ & temporary file's path.\\
    $Y$ & whether key caching is enabled.\\
    $L$ & number of memory lanes for concurrency.\\
    $T$ & number of tasks per lane segment.\\
    $H$ & number of lane segments per hard-disk read.\\
    $B$ & minimum \emph{caveman's entropy bits} to inject into $p$.\\
    $K$ & output key's size in bytes.\\
\end{tabularx}

\subsubsection{internal variables}
\begin{tabularx}{\columnwidth}{lX}
    $\enc$      & encryption function.\\
    $\hash$     & hashing function.\\
    $\fread$    & hard-disk file reading function, with seeking.  e.g.
                    $\fread(x, y, z)$ reads $z$ many bytes from file $x$
                    after seeking $y$ bytes forward.\\
    $\fwrite$   & hard-disk file writing function.\\
    $C$         & $\gets \begin{cases}
                        64 \text{ bytes} & \text{if $\enc$ is
                                            \emph{xchacha20}}\\
                        16 \text{ bytes} & \text{if $\enc$ is \emph{aes}}\\
                        \vdots & \\
                     \end{cases}$\\
                & this to reflect the block size of the encryption
                    algorithm that implements $\enc$.\\
\end{tabularx}
\begin{tabularx}{\columnwidth}{lX}
    $V$ & $\gets \begin{cases}
                        32 \text{ bytes} & \text{if $\enc$ is
                                            \emph{xchacha20}}\\
                        16 \text{ bytes} & \text{if $\enc$ is
                            \emph{aes-128}}\\
                        32 \text{ bytes} & \text{if $\enc$ is
                            \emph{aes-256}}\\
                        \vdots & \\
                     \end{cases}$\\
                & this is the size of the encryption key that's used to
                    solve \emph{ciphart}'s tasks.  this is different than
                    output key's size, $K$, which is $\enc$-independent.\\
    $\hat T$    & $\gets \maxf(\lceil V C^{-1}\rceil, T)$.  this
                    is to ensure that we have enough encrypted bytes for
                    new keys.\\
    $\hat T$    & $\gets \hat T - (\hat T \bmod 2) + 2$.  this is to ensure
                    that there is an even number of tasks in a segment.
                    why?  because we need a buffer for storing the
                    clear-text and another for storing the output
                    cipher-text.\\
    $\hat M$    & $\gets M - (M \bmod C\hat TL) + C\hat TL$.  this is to
                    ensure that it is in multiples of $C\hat TL$.  why?  so
                    that all segments are of equal lengths in order to
                    simplify \emph{ciphart}'s logic.  e.g. it wouldn't be
                    nice if the last segments were of unequal sizes.\\
    $G$         & $\gets \hat MC^{-1}\hat T^{-1}L^{-1}$.  total number of
                    segments per lane.\\
    $n$    & $\gets 0$.  actual number of times $\enc$ is called.\\
    $m_i$       & $C$-bytes memory for $i^{th}$ task in the $\hat M$-bytes
                    pad.\\
    $n_l$       & $\gets \maxf(\text{nonce})L^{-1}l$.  nonce variable for
                    $l^{th}$ lane.\\
    $f$         & $\gets 0$.  a counter indicating number of times
                    memory is filled with $\hat M$ many bytes.\\
    $d$         & $\gets 0$.  a counter indicating total number of saved
                    blocks into hard-disk.\\
    $h$         & $\gets 0$.  a counter indicating number of processed lane
                    segments since the last hard-disk read.\\
    $u$         & $\gets 0$.  a counter indicating number of times key was
                    updated from the hard-disk.\\
    $*v$         & $\gets *\hash(P \Vert S \Vert M \Vert D \Vert \ldots
                    \Vert K, V)$.  a pointer to the first byte where
                    $V$-bytes key is stored.  $v$ is the key itself, and
                    $*v$ is a pointer to it.\\
    $Z$         & $\gets \hash(v \Vert 0, V)$.  file name where output key,
                    $k$, is expected to be cached, if $k$ was previously
                    cached.\\
\end{tabularx}

\subsubsection{output}
\begin{tabularx}{\columnwidth}{lX}
$k$ & $K$-bytes key.\\
$n$ & total number of times $\enc$ was actually called.  $\log_2 n$ is
        total number shannon's entropy bits that \emph{ciphart} injected
        into $k$, such that $\log_2 n \ge B$.\\
$\hat M$ & actual number of bytes required to exist in random-access
            memory.\\
$d$     & actual number of bytes required to exist in hard-disk.\\
\end{tabularx}

\subsubsection{steps}
steps of \emph{ciphart} is shown in \cref{alg_ciphart}.  this corresponds
to \emph{argon2d}.  adding a \emph{ciphart-i} variant is a trivial matter,
i just didn't do it yet because my threat model currently doesn't benefit
from a password independent variant.

\begin{algorithm}
\If{$Y$ {\bf and} $\fexists(Z)$}{\label{ciphart_cache_hit}
    $k, n, d \gets \fread(Z, 0, V + \sizeof(k\Vert n\Vert d))$\;
    \Return{$k, n, d$}
}
\While{$1$}{
    \For{$g \gets 0, 1, \ldots, G-1$}{
        \For{$l \gets 0, 1, \ldots, L-1$}{\label{ciphart_lanes}
            \For{$t \gets 0, 1, \ldots, T-1$}{
                $i \gets gLT + lT + t$\;
                \uIf{$t < T - 1$}{
                    $j \gets i + 1$\;
                }\ElseIf{$t = T - 1$}{
                    $j \gets i - T + 1$\;
                }
                $m_j \gets \enc(m_i, n_l, *v)$\;
                $n_l \gets n_l + 1$\;
                \uIf{$f = 0$}{
                    $*v \gets m_j \bmod (gLTC + tC - V)$\;
                    \If{$*v \ge gLTC - V$}{
                        $*v \gets *v + lTC$\;
                    }
                }\Else{
                    $*v \gets m_j \bmod (\hat M - LTC + tC - V)$\;
                    \uIf{$*v \ge gLTC + tC - V$}{
                        $*v \gets *v + LTC$\;
                    }\ElseIf{$*v \ge gLTC - V$}{
                        $*v \gets *v + lTC$\;
                    }
                }
            }
        }
        $n \gets n + LT$\;
        \uIf{$d \le D$}{\label{ciphart_D_filling}
            \For{$i \gets gLT, \ldots, gLT+(T-1)$}{
                $\fwrite(F, m_i)$\;
                $d \gets d + 1$\;
                \If{$d \ge D$}{ \textbf{break}\; }
            }
        }\Else{
            $h \gets h + L$\;
            \If{$h \ge H$}{\label{ciphart_D_using}
                $*v \gets *\fread(F, v \bmod (d - V), V)$\;
                $u \gets u + 1$\;
                $h \gets 0$\;
            }
            \If{$f \ge 1$ \bf and $u \ge 1$ \bf and $n \ge 2^B$}{
                $g_{\text{last}} \gets g$\;
                \textbf{go to} \cref{ciphart_out}\;
            }
        }
    }
    $f \gets f+1$\;
}
$i \gets g_{\text{last}}LT$\;\label{ciphart_out}
$k \gets \hash(m_{i+0T} \Vert m_{i+1T} \Vert \ldots \Vert m_{i+(L-1)T}, K)$\;
\If{$Y$}{\label{ciphart_cache_miss}
    $\fwrite(Z, k \Vert n \Vert d)$\;
}
$\delete(F)$\;
\Return{$k, n, d$}
\caption{ciphart}
\label{alg_ciphart}
\end{algorithm}

\subsection{noteworthy features}
\subsubsection{parallelism}
since iterations of the loop in \cref{ciphart_lanes} in \cref{alg_ciphart}
are fully independent of one other, they can quite happily utilise $L$ cpu
cores, specially when segment sizes, $T$, are larger.

\subsubsection{memory-hardness}
\begin{proof}
    \cref{alg_ciphart} is just a variation of \emph{argon2d}.  so if
    \emph{argon2d} is memory-hard, then so is \emph{ciphart}.
\end{proof}


\subsubsection{memory-\emph{harder}ness}
thanks to discovery \ref{discov_key_caching}, we can cache keys, as done in
\cref{ciphart_cache_hit}, without increasing assumptions of the threat
model of the vast majority of users.  then, memory-\emph{harder}ness
becomes possible.  this process goes like this:
\begin{enumerate}
    \item starts by running an $\enc$-based variant of \emph{argon2},
    except that, as it is going, it keeps writing the updated segments into
    the hard-disk until the size of it satisfies the $D$ bytes limit.  this
    is shown in \cref{ciphart_D_filling} in \cref{alg_ciphart}.

    optimising this hard-disk filling with $D$ bytes is not a big deal,
    since this feature is probably going to be used only once; thanks to
    key caching.  that said:
    \begin{itemize}
        \item this feature is optional.  i.e. in case someone doesn't like
        the hard-disk caching, he can set $D \gets 0$ to disable it.  but,
        for most people, i don't understand why you would want to disable
        it.  e.g. if you're already trusting \texttt{root}, then i think
        that you can use this feature without changing your threat model.

        i personally like it a lot as it allows me to achieve
        memory-\emph{harder}ness way beyond my random-access memory.  just
        imagine the look on the face of those asics crackers once they hear
        that your \emph{ciphart} requires, say, $50$ giga bytes!

        \item when key caching is enabled, i.e. $Y \gets 1$, this hard-disk
        writing is done only initially, and subsequent uses appear as
        almost instantaneous.

        \item this hard-disk writing can be slightly optimised by using a
        non-blocking write operation.  but i think that the trivial
        reduction in this time may not justify increasing  code's
        complexity, so i don't plan to implement non-blocking writes in
        \texttt{libciphart}.
    \end{itemize}

    \item then, once the $D$ bytes hard-disk requirement is satisfied, the
    process continues to run the modified $\enc$-based variant of
    \emph{argon2} except that it updates the key $v$ from those $D$ bytes
    every time $S$ many segments are solved.  this step makes
    \emph{ciphart} require $D+M$ bytes.  this is shown in
    \cref{ciphart_D_using} in \cref{alg_ciphart}.

    if $S$ is large enough, this can be done efficiently without blocking
    the cpu noticeably, as the randomly obtained $V$ bytes can be read by
    using the $O(1)$ operation, \texttt{seek}, over the $D$ bytes.

    \item delete the $D$ bytes, just to free up the disk space.  no need to
    securely delete those $D$ bytes since we can save them using a
    temporary random key that's forgotten later on.
\end{enumerate}


\subsection{comparison}

\section{application scenarios}
\subsection{a currently-useful scenario}
user has a password manager which generates unique $256$ bit entropy keys for
each online service that he uses.  the user also renews keys for his online
services every now and then.   so his online accounts generally have high
security.

but user's problem is how to lock and unlock his password manager's
passwords database.  should he use a physical usb-stick key that types a
high-entropy key that encrypts and decrypts the database?   the user
doesn't want this physical key because of several reasons:
\begin{itemize}
    \item he tends to lose his keys a lot, and, for certain tasks, the risk
    of needing to wait for until he gets a backup usb-stick key is too
    much.

    \item he doesn't want to be caught having cryptographic usb-stick keys,
    because as such is an evidence that he has encrypted content.  the user
    wants to have the choice to lie that he has no clue.  so not having
    usb-stick keys with him helps his case to lie.

    \item he wants to be torture-resistant so that an adversary cannot
    forcefully take his keys from him in order to login into his services.
    he may rather want to die than to give the password to the adversary.
    this works because, so far, the brain is a pretty private information
    store.
\end{itemize}

in this scenario, the user memorises a sensible password that he can
remember, with enough initial entropy, and then uses \emph{ciphart},
preferably with disk caching, to inject large amounts of entropy bits into
his derived keys, way beyond the reach of pre-\emph{ciphart} key derivation
functions; thanks to \cref{theorem_entropy_injection} and discovery
\ref{discov_key_caching}.

here, the expensively derived key is only used to unlock a local password
manager, which offers a protection against situations where a backup copy
of the passwords database is stolen.  this may enable the user to store his
passwords manager in an online file synchronisation service for more
convenient system migrations to further reduce his login delays should he
face the need to migrate to a new, say, personal computer.

\begin{advertisement}
    in case you're interested in such a passwords manager, i've also made
    \emph{nsapass}\footnote{\url{https://github.com/Al-Caveman/nsapass}}
    --- a flexible and a simple passwords manager in a few hundreds of
    python lines of code, that uses \emph{ciphart} by default.  i think
    this is the best command-line interfacing passwords manager by far, for
    its usability, and for the fact that auditing it is easy, thanks for it
    being only in a few hundreds of python lines.
\end{advertisement}

\subsection{a later-useful scenario}
all input password fields will internally call \texttt{libciphart} to
derive more expensive keys.  this way, applications, such as
\emph{firefox}, \emph{mutt}, \ldots, will never send actual passwords, but
will only send \emph{ciphart}-derived keys with increased shannon's entropy
content.

thanks to \cref{theorem_entropy_injection}, this will automatically
increase shannon's entropy of all users' passwords without requiring users
to memorise harder passwords.  thanks to discovery
\ref{discov_key_caching}, the user will not face any delay in his daily
use, except only an initial delay to create the cache entry.

i also think that it is \emph{generally} better to have expensive key
derivation functions in the client side as opposed to the server side.
because remote servers always have the incentive to reduce the complexity
of the key derivation function in order to free more resources for other
things that bring them money.

perhaps, one day in the future, the \emph{password} field in \emph{html}
should have an argument called \emph{kdf} that goes like \texttt{<input
type="password" kdf="ciphart" \ldots>} which specifies the key derivation
function that the web browser must use to derive a more secure key, and
then send the derived key instead of sending the password.  browsers are
already kind enough to \emph{manage} our passwords for us, so this will be
just a matter of making them send more secure key as opposed to the less
secure passwords.  

i personally don't let browsers manage my passwords, because i don't trust
them, and because i don't need them, as i use my own passwords manager
\emph{nsapass}.  so i don't care about this.  i just wrote it in case it
helps the cute normal people more securely syndicate videos of their happy
puppies across the interwebs, without requiring them to become better
people.  why?  i guess i'm too kind, and i think these people are sort of
adorable the way they are, so i don't mind to preserve their innocence.

either  way, \emph{ciphart} is perfectly usable on the server side.  it is
just that i think \emph{ciphart}, \emph{argon2}, \emph{scrypt}, \ldots make
better sense when placed on the client side.  ideally implemented in a
passwords manager like \emph{nsapass} as i said earlier, or integrated into
the password fields of web browsers in order to preserve the innocence of
the adorable normies.

\vfill
\break
\appendix
\section{donations}
this work is sponsored by an unexpected nerdiness.  nothing here was
supposed to happen.  the cause  remains unknown, so we call it
\emph{something random with a large entropy}.

not-so-ancient mythology\footnote{i made this myth up.} has it that, every
time a donation is made for a good cause, a beautiful torch is kindled
somewhere deep within an otherwise cold and dark space.

the torch is pure with an innocent smile every time a spending is committed
towards good.  but, she drops a tear when the spending is done in excess.
despite her attempts to maintain a steady light, sometimes her tears force
her light to flicker, and sometimes the flickering causes her to disappear.
when she is gone, she is never back again, except in the memories of those
who have once witnessed her charm\ldots

so, if you opt for a donation, i'd appreciate exercising wisdom, so that
the sensible balance is not exceeded; lest we want a beautiful torch to be
born, only to make her disappear in her tears.

\subsection{bitcoin}
\begin{center}
    \includegraphics[width=121px]{./pics/btc_wallet_address_trimmed.png}
    \texttt{bc1qtylzjtgd0yu4v7f8hyfzufn7nu692v9fc88jln}
\end{center}

\end{document}
